% ============================================================================
% Fichier  : partie4/P04_C17_forensique_cloud.tex
% Titre    : Forensique Cloud
% Auteur   : MINKA MI NGUIDJOÏ Thierry Emmanuel
% Version  : 1.0 -- Juin 2026
% Statut   : RÉDIGÉ
% Sources  : Ch. 10 (portails LEA, art. 29 Budapest, Protocole II)
%            Ch. 11 (CLOUD Act, SCA, RGPD art. 23)
%            Ch. 15 (logs, corrélation, CloudTrail mentionné)
%            Affaires WOURI (Gmail), BAOBAB (Binance), MVOM (AWS)
% Positionnement : Ch. 10-11 ont traité le cadre légal de l'accès aux
%                  données cloud (portails LEA, MLAT, CLOUD Act). Ce chapitre
%                  traite le "comment" technique : que demander, comment
%                  analyser les logs obtenus, comment structurer les preuves
%                  cloud pour le tribunal.
% ============================================================================

\chapter{Forensique Cloud}
\label{chap:for-cloud}

\epigraph{%
    « The cloud is just someone else's computer. »%
}{--- Adage des administrateurs système}

\niveauChapitre{Expert}{%
    Chapitres~\ref{chap:norm-glob} (portails LEA, Budapest),
    \ref{chap:leg-mond} (CLOUD Act, SCA, RGPD) et
    \ref{chap:for-res} (logs, corrélation) indispensables.%
}

\begin{boxobjectifs}
\begin{itemize}
    \item Comprendre les modèles de service cloud (SaaS, PaaS, IaaS)
          et leur impact sur la disponibilité des preuves forensiques.
    \item Identifier les logs disponibles sur les trois grandes
          plateformes (AWS, Azure, Google Cloud) et leur contenu
          forensique.
    \item Formuler des demandes structurées aux portails LEA des
          grands fournisseurs~; obtenir et analyser les réponses.
    \item Analyser les journaux AWS CloudTrail et Azure Activity Log
          pour reconstituer une chronologie d'attaque cloud.
    \item Gérer la tension entre éphéméralité des données cloud
          et exigences probatoires de persistance.
    \item Coordonner une investigation cloud impliquant plusieurs
          juridictions et fournisseurs.
\end{itemize}
\end{boxobjectifs}

\bigskip

% ============================================================================
\section*{Le problème central de la forensique cloud}
% ============================================================================

La forensique cloud pose trois défis que les investigations classiques
ne posent pas~:

\begin{enumerate}
    \item \textbf{Souveraineté diffuse}~: les données sont stockées dans
          des datacenters dont la localisation peut changer dynamiquement
          et dépend du fournisseur, pas de l'enquêteur. Qui a juridiction~?
    \item \textbf{Éphéméralité}~: les instances cloud peuvent être
          détruites en secondes, emportant toutes leurs données. La fenêtre
          d'acquisition est souvent de minutes, pas d'heures.
    \item \textbf{Multi-tenancy}~: les données de plusieurs clients
          coexistent sur la même infrastructure physique~; l'isolement
          est logique, pas physique.
\end{enumerate}

\separateur

% ============================================================================
\section{Modèles de Service et Disponibilité des Preuves}
\label{sec:modeles-cloud}
% ============================================================================

Le modèle de service détermine qui contrôle quoi --- et donc qui peut
fournir quelles preuves à la demande de l'investigateur.

\begin{tableaucours}{p{2.5cm}p{3.0cm}p{3.0cm}p{3.5cm}}{%
    \tableauHeader{Modèle} &
    \tableauHeader{Ce que gère le client} &
    \tableauHeader{Ce que gère le fournisseur} &
    \tableauHeader{Preuves disponibles via LEA}
}
\textbf{SaaS} (Gmail, Office~365, Dropbox) &
    Ses données uniquement &
    Tout (infra, OS, appli, logs) &
    Logs de connexion, métadonnées, contenu (MLAT/warrant) \\
\textbf{PaaS} (Heroku, Google App Engine) &
    Code applicatif, données &
    Infra, OS, runtime &
    Logs plateforme + logs applicatifs si conservés \\
\textbf{IaaS} (AWS EC2, Azure VM, GCP Compute) &
    OS, appli, données, logs &
    Physique, hyperviseur &
    Logs de la couche cloud uniquement~; logs OS~= réquisition
    au client (ou saisie du compte) \\
\end{tableaucours}
\vspace{-0.8em}
\begin{center}{\small\itshape Tableau~17.1 --- Modèles cloud et disponibilité des preuves}\end{center}

\begin{boiteRemarque}{Implication pratique~: SaaS est plus simple que IaaS}
Paradoxalement, l'investigation d'un compte Gmail (SaaS) est plus
directe qu'une instance AWS compromise (IaaS). Pour Gmail, une demande
au portail LEA Google suffit pour les métadonnées~; pour une instance
EC2, le fournisseur ne voit pas le contenu des VMs~--- il faut soit
accéder au compte AWS compromis, soit saisir légalement l'instance.
\end{boiteRemarque}

% ============================================================================
\section{Logs Cloud~: ce qui Existe et où le Trouver}
\label{sec:logs-cloud}
% ============================================================================

\subsection{AWS CloudTrail~: l'audit complet des API}

\termecle{AWS CloudTrail}\index{CloudTrail} enregistre chaque appel API
effectué sur un compte AWS~: qui a fait quoi, quand, depuis quelle IP.
C'est le journal d'audit le plus complet du cloud AWS.

\begin{boxoutils}{Analyse AWS CloudTrail}
\begin{lstlisting}[language=bash_forensique]
# CloudTrail stocke ses logs dans S3 au format JSON (gzip)
# Structure d'un log CloudTrail type :
# {
#   "eventTime": "2026-03-13T02:15:33Z",
#   "eventSource": "ec2.amazonaws.com",
#   "eventName": "DescribeInstances",
#   "userIdentity": {
#     "type": "IAMUser",
#     "userName": "admin",
#     "arn": "arn:aws:iam::123456789:user/admin"
#   },
#   "sourceIPAddress": "185.220.XX.XX",
#   "userAgent": "aws-cli/2.11.0 Python/3.11.0"
# }

# Telecharger et analyser les logs CloudTrail avec AWS CLI
aws s3 ls s3://mon-bucket-cloudtrail/AWSLogs/123456789/
aws s3 sync s3://mon-bucket-cloudtrail/ /output/cloudtrail/

# Analyser avec jq (JSON query)
# Toutes les actions depuis une IP suspecte
cat /output/cloudtrail/*.json.gz | gunzip | \
    jq -r 'select(.sourceIPAddress == "185.220.XX.XX") |
    [.eventTime, .eventName, .userIdentity.userName,
     .sourceIPAddress] | @csv'

# Actions de creation/suppression d'instances (indicateurs)
cat /output/cloudtrail/*.json.gz | gunzip | \
    jq -r 'select(.eventName | test("RunInstances|TerminateInstances|
        CreateUser|DeleteUser|AttachUserPolicy")) |
    [.eventTime, .eventName, .userIdentity.userName,
     .sourceIPAddress] | @csv'

# Verifier si CloudTrail a ete desactive (anti-forensique cloud)
cat /output/cloudtrail/*.json.gz | gunzip | \
    jq -r 'select(.eventName ==
        "StopLogging" or .eventName == "DeleteTrail") |
    [.eventTime, .eventName, .userIdentity.userName] | @csv'
# Trouver "StopLogging" dans les logs -> tentative d'anti-forensique
\end{lstlisting}
\end{boxoutils}

\begin{boxscenario}{CloudTrail --- reconstitution d'une compromission AWS MVOM}
L'analyse des logs CloudTrail du compte AWS de TRANSBOIS~SA révèle
la séquence suivante~:

\begin{lstlisting}[language=bash_forensique]
# Timeline CloudTrail filtree sur 13/03/2026 02h00-03h00
02:13:45  AssumeRole       IP:185.220.XX.XX  Role:EC2-Admin
02:14:03  DescribeInstances IP:185.220.XX.XX  User:admin
02:14:11  GetConsoleOutput  IP:185.220.XX.XX  Instance:i-0a1b2c3d4e
02:16:02  RunInstances      IP:185.220.XX.XX  Type:t2.xlarge
          # Lancement d'une instance suspecte
02:17:55  PutObject         IP:185.220.XX.XX  Bucket:transbois-backup
          # Exfiltration vers S3 (45 Mo)
02:19:48  StopLogging       IP:185.220.XX.XX
          # TENTATIVE SUPPRESSION DES LOGS CloudTrail
02:19:49  [CloudTrail arrêté -- fenêtre aveugle]
02:23:11  TerminateInstances  # Instance suspecte detruite
\end{lstlisting}

\textbf{Constatation~:} L'action \texttt{StopLogging} à 02h19:49 a
interrompu la journalisation CloudTrail. Cependant, les logs antérieurs
(02h13--02h19) sont déjà enregistrés dans S3 et ne peuvent pas être
rétroactivement effacés si le bucket S3 est protégé
(\textit{Object Lock / WORM})~--- ce qui est le cas ici.
\textbf{L'anti-forensique cloud a échoué car les logs étaient
en écriture unique.}
\end{boxscenario}

\subsection{Azure Activity Log et Azure AD Sign-In Logs}

\begin{boxoutils}{Analyse Azure Activity Log et Sign-In}
\begin{lstlisting}[language=bash_forensique]
# Azure CLI -- extraire les logs d'activite
az monitor activity-log list \
    --start-time "2026-03-13T02:00:00Z" \
    --end-time "2026-03-13T03:00:00Z" \
    --output json > /output/azure_activity.json

# Analyser avec jq
cat /output/azure_activity.json | jq -r \
    '.[] | [.eventTimestamp, .operationName.value,
    .caller, .properties.statusCode] | @csv'

# Azure AD Sign-In Logs (tentatives de connexion)
# Via Azure Portal -> Azure Active Directory -> Sign-ins
# Ou via Microsoft Graph API :
curl -H "Authorization: Bearer $TOKEN" \
    "https://graph.microsoft.com/v1.0/auditLogs/signIns?\
    \$filter=createdDateTime ge 2026-03-13T02:00:00Z and
    createdDateTime le 2026-03-13T03:00:00Z" \
    > /output/azure_signins.json

cat /output/azure_signins.json | jq -r \
    '.value[] | [.createdDateTime, .userDisplayName,
    .ipAddress, .status.errorCode, .location.city] | @csv'
# errorCode 0 = succes, 50126 = mot de passe incorrect, etc.
\end{lstlisting}
\end{boxoutils}

\subsection{Google Cloud Logging~: Cloud Audit Logs}

\begin{lstlisting}[language=bash_forensique,
    caption={Google Cloud Audit Logs --- analyse forensique}]
# Google Cloud Logging via gcloud CLI
gcloud logging read \
    "timestamp >= \"2026-03-13T02:00:00Z\"
     AND timestamp <= \"2026-03-13T03:00:00Z\"
     AND protoPayload.methodName:(
        \"compute.instances.insert\"
        OR \"compute.instances.delete\"
        OR \"storage.buckets.create\")
     AND protoPayload.authenticationInfo.principalEmail != \"\"" \
    --project=mon-projet \
    --format=json > /output/gcp_audit.json

# Extraire les createurs de ressources (qui a lance quoi)
cat /output/gcp_audit.json | python3 -c "
import json, sys
for entry in json.load(sys.stdin):
    pp = entry.get('protoPayload', {})
    print(
        entry.get('timestamp',''),
        pp.get('methodName',''),
        pp.get('authenticationInfo',{}).get('principalEmail',''),
        pp.get('requestMetadata',{}).get('callerIp','')
    )
"
\end{lstlisting}

% ============================================================================
\section{Obtenir des Preuves Auprès des Fournisseurs SaaS}
\label{sec:preuve-saas}
% ============================================================================

\subsection{Structurer une demande au portail LEA}

Le Chapitre~\ref{chap:norm-glob} (Section~\ref{sec:mlat-lea}) a
présenté les portails LEA. Ce chapitre approfondit la \emph{structure}
de la demande elle-même~--- les éléments sans lesquels elle sera
systématiquement rejetée.

\begin{tableaucours}{p{3.5cm}p{8.0cm}}{%
    \tableauHeader{Champ obligatoire} &
    \tableauHeader{Contenu requis et formulation recommandée}
}
\textbf{Type de demande} &
    Preservation Request (urgence)~; Production Request (standard)~;
    Emergency Request (menace imminente). Choisir le bon type~:
    une Production sans Preservation préalable arrive souvent après
    que les données ont été effacées. \\
\textbf{Identifiants du compte} &
    Adresse email, numéro de téléphone, URL du profil, IMEI associé.
    Toujours donner au moins deux identifiants croisés pour lever
    l'ambiguïté. \\
\textbf{Base légale} &
    Citer explicitement~: art.~52 et~57 \loicam{2010/012}~;
    Budapest art.~16--17 si l'État requis est partie~;
    SCA 18~U.S.C.~\S~2703 pour fournisseurs US. \\
\textbf{Qualification de l'infraction} &
    Description factuelle + qualification juridique des deux pays
    (double incrimination). Pour un fournisseur US~: citer le CFAA. \\
\textbf{Période visée} &
    Date de début et de fin précises. Éviter les demandes ouvertes
    (``toute la vie du compte'')~: systématiquement rejetées. \\
\textbf{Données demandées} &
    Distinguer métadonnées (abonné, connexion) et contenu~:
    les métadonnées s'obtiennent sans warrant~; le contenu nécessite
    un MLAT ou équivalent. \\
\textbf{Coordonnées de contact} &
    Nom, titre, email officiel, numéro de téléphone direct.
    Les réponses urgentes vont directement à ce contact. \\
\end{tableaucours}
\vspace{-0.8em}
\begin{center}{\small\itshape Tableau~17.2 --- Éléments obligatoires d'une demande LEA}\end{center}

\subsection{Analyser la réponse d'un fournisseur}

\begin{boxoutils}{Traitement d'une réponse Google (métadonnées Gmail)}
\begin{lstlisting}[language=bash_forensique]
# Structure typique d'une reponse Google aux autorites
# (format JSON ou fichier texte structure)
# Exemple : reponse pour un compte Gmail suspect

# 1. Informations d'abonne (disponibles sans warrant)
{
  "accountEmail": "suspect.wouri@gmail.com",
  "googleId": "123456789012345678901",
  "accountCreated": "2025-12-28T10:15:22Z",
  "lastLoginDate": "2026-01-15T14:35:11Z",
  "recoveryEmail": "autre.compte@yahoo.fr",
  "recoveryPhone": "+23365XXXXXXX",
  "services": ["Gmail", "Google Drive", "Google Photos"]
}

# 2. Logs de connexion (disponibles sans warrant)
[
  {"loginTime": "2026-01-15T14:28:01Z",
   "ipAddress": "41.82.XXX.XXX",       # IP senegale (Sonatel)
   "loginType": "GOOGLE_LOGIN",
   "success": true},
  {"loginTime": "2026-01-15T14:35:09Z",
   "ipAddress": "41.82.XXX.XXX",
   "loginType": "GMAIL_SMTP_AUTH",
   "success": true}
]

# Croiser avec la capture reseau WOURI (Ch. 15) :
# Email SMTP envoye depuis contact@medequip-portugal.com a 14h28:03
# Connexion Gmail depuis IP senegalaise a 14h28:01
# -> le compte suspect s'est connecte 2 secondes avant l'envoi
\end{lstlisting}
\end{boxoutils}

\begin{boxjuris}
\textbf{Valeur probatoire d'une réponse de fournisseur}

Une réponse d'un fournisseur US (Google, Meta, Microsoft) à une
demande LEA constitue un \textit{business record} au sens de la
FRE~Rule~803(6) (Chapitre~\ref{chap:leg-mond}). Elle est directement
admissible en droit américain. En droit camerounais~:
\begin{itemize}
    \item Considérée comme un document étranger authentifié~;
    \item Nécessite l'attestation d'un expert judiciaire expliquant
          le processus de génération et la fiabilité du système~;
    \item Hash SHA-256 du fichier reçu à calculer et documenter
          immédiatement (custody numérique des réponses LEA,
          cf. Ch.~\ref{chap:custody}).
\end{itemize}
\end{boxjuris}

% ============================================================================
\section{Gérer l'Éphéméralité~: Conservation d'Urgence}
\label{sec:ephemerality}
% ============================================================================

\subsection{Le problème de la rétention limitée}

Les données cloud ont des durées de vie variables selon les fournisseurs~:

\begin{tableaucours}{p{3.5cm}p{3.5cm}p{4.5cm}}{%
    \tableauHeader{Type de donnée} &
    \tableauHeader{Durée de rétention typique} &
    \tableauHeader{Action à entreprendre}
}
Logs de connexion IP & 30--90~jours selon fournisseur &
    Demande de conservation \textbf{immédiate} via art.~29 Budapest
    ou portail LEA \\
Contenu emails supprimés & 30~jours (corbeille) puis effacé &
    Preservation Request \textbf{avant} la fin des 30~jours \\
Instances cloud éteintes & Selon configuration (souvent~: 7~jours) &
    Saisie immédiate ou snapshot forensique \\
Données Snap/Telegram & Heures à jours &
    Urgence absolue~; canal INTERPOL 24/7 \\
Transactions financières & Souvent 5--10~ans (réglementation AML) &
    Standard MLAT acceptable \\
\end{tableaucours}
\vspace{-0.8em}
\begin{center}{\small\itshape Tableau~17.3 --- Durées de rétention et urgence d'action}\end{center}

\subsection{Snapshot forensique d'une instance cloud}

Quand l'investigateur a accès au compte cloud compromis (ou que le
suspect est le propriétaire du compte et coopère), le snapshot est
l'équivalent cloud de l'image E01~:

\begin{boxoutils}{Snapshot forensique AWS / GCP}
\begin{lstlisting}[language=bash_forensique]
# AWS -- snapshot d'un volume EBS (disque de VM)
# PREREQUIS : acces au compte AWS (credentials)

# 1. Identifier le volume de l'instance
aws ec2 describe-instances --instance-ids i-0a1b2c3d4e \
    --query 'Reservations[].Instances[].BlockDeviceMappings[].Ebs.VolumeId'
# Resultat : vol-0xyz789abc

# 2. Creer un snapshot (image read-only du volume)
aws ec2 create-snapshot --volume-id vol-0xyz789abc \
    --description "Forensic snapshot MVOM 2026-03-13 OPJ NKENG"
# Resultat : snap-0forensic123

# 3. Copier le snapshot dans un compte forensique isole
aws ec2 copy-snapshot --source-region eu-west-1 \
    --source-snapshot-id snap-0forensic123 \
    --region eu-west-1 \
    --description "Forensic copy - case DOS-2026-034"

# 4. Calculer et documenter le hash du snapshot
# (via AWS S3 Export puis hash local)
aws ec2 export-image --snapshot-ids snap-0forensic123 \
    --s3-export-location S3Bucket=forensic-exports,S3Prefix=case-034/
# Apres export :
aws s3 cp s3://forensic-exports/case-034/snapshot.vmdk /output/
sha256sum /output/snapshot.vmdk > /output/snapshot_hash.txt

# GCP -- snapshot d'un disque Persistent Disk
gcloud compute disks snapshot DISK_NAME \
    --snapshot-names=forensic-snap-20260313 \
    --zone=europe-west1-b \
    --description="Forensic snapshot MVOM DOS-2026-034"
\end{lstlisting}
\end{boxoutils}

\begin{boiteRemarque}{Le snapshot n'est pas l'image E01}
Un snapshot cloud est une copie logique cohérente du volume à un instant t.
Il préserve le contenu mais \textbf{pas exactement les mêmes garanties
qu'un write blocker physique}~: la création du snapshot est elle-même
une opération enregistrée dans CloudTrail. Documenter obligatoirement~:
\begin{enumerate}
    \item Horodatage exact de la création du snapshot~;
    \item ID du snapshot~;
    \item Compte AWS/GCP utilisé pour le créer~;
    \item Hash SHA-256 du fichier exporté.
\end{enumerate}
\end{boiteRemarque}

% ============================================================================
\section{Investigation Multi-Fournisseurs~: Corrélation Cloud}
\label{sec:correlation-cloud}
% ============================================================================

Les affaires complexes impliquent souvent plusieurs fournisseurs cloud
simultanément. La corrélation entre leurs données est la clé pour
reconstituer l'image complète.

\begin{boxscenario}{Affaire BAOBAB --- corrélation multi-cloud}
L'Opération BAOBAB (fraudes MoMo, Chapitre~\ref{chap:norm-glob}) implique~:
\begin{itemize}
    \item Un compte Binance (exchange crypto) via lequel les fonds
          ont été convertis~;
    \item Une adresse email Yahoo utilisée pour créer le compte Binance~;
    \item Un VPN hébergé sur AWS EC2 (masquage de l'IP réelle)~;
    \item Un wallet Bitcoin avec transactions sur la blockchain.
\end{itemize}

\textbf{Séquence de demandes parallèles~:}

\begin{enumerate}
    \item \textbf{Preservation simultanée~:} portail LEA Binance +
          portail LEA Yahoo + AWS CloudTrail. \textbf{Simultané,
          pas séquentiel}~: les données se recoupent et les effacer
          l'une révèlerait l'enquête à l'autre.

    \item \textbf{Analyse blockchain (immédiate, sans demande)~:}\\
          \texttt{https://www.blockchain.com/btc/address/1BaobabXXX}\\
          L'adresse Bitcoin et toutes ses transactions sont publiques.
          Identifier les échanges de destination (exchanges~: ont des
          obligations KYC) et les wallets liés.

    \item \textbf{Corrélation temporelle~:}
\begin{lstlisting}[language=bash_forensique]
# Chronologie fusionnee des 3 fournisseurs :
# [MoMo SMS]   2026-01-15 14:22:45  Transaction 250000 FCFA
# [Yahoo LEA]  2026-01-15 14:23:12  Login compte Yahoo IP:41.82.XXX.XXX
# [Binance LEA] 2026-01-15 14:23:45 Login Binance IP:41.82.XXX.XXX
# [Binance LEA] 2026-01-15 14:24:02 Conversion XAF -> USDT (240000 FCFA)
# [AWS CT]     2026-01-15 14:25:11  Connection SSH to EC2 VPN
# [Blockchain] 2026-01-15 14:31:05  Transfer 185 USDT -> Bitcoin wallet
\end{lstlisting}
          La même IP sénégalaise (\texttt{41.82.XXX.XXX}) apparaît
          dans les logs Yahoo et Binance à 78~secondes d'intervalle~---
          cohérence forte.
    \end{enumerate}
\end{boxscenario}

\begin{boxCRO}
\textbf{Analyse CRO --- Blockchain comme preuve forensique}

$\mathcal{R}$ (Fiabilité)~: la blockchain Bitcoin est immuable par
conception. Une transaction enregistrée est permanente et vérifiable
indépendamment par tout analyste avec les mêmes outils. La fiabilité
est maximale~--- bien supérieure à un log de serveur qu'un administrateur
pourrait modifier.

$\mathcal{O}$ (Opposabilité)~: l'immuabilité de la blockchain constitue
une preuve d'une opposabilité exceptionnelle~; elle ne peut pas être
contestée sur le fond (la transaction a eu lieu). Elle peut être
contestée sur l'attribution (``prouvez que je contrôlais ce wallet'').

$\mathcal{C}$ (Confidentialité)~: paradoxalement, la blockchain est
\textbf{publique et pseudonyme}. L'anonymat peut être levé en
identifiant les échanges KYC aux deux extrémités de la chaîne.

\cro{$\uparrow$ pseudonyme mais levable}{$\uparrow\uparrow\uparrow$ immuabilité maximale}{$\uparrow\uparrow\uparrow$ incontestable}
\end{boxCRO}

% ============================================================================
\section{Défis Spécifiques~: Chiffrement et Résidualité}
\label{sec:defis-cloud}
% ============================================================================

\subsection{Données chiffrées côté client}

Certains services (Proton Mail, iCloud avec E2E, Signal) chiffrent les
données \textbf{côté client} avant de les envoyer au serveur~: même
avec un warrant, le fournisseur ne peut pas déchiffrer.

\begin{tableaucours}{p{3.5cm}p{3.0cm}p{5.0cm}}{%
    \tableauHeader{Service} &
    \tableauHeader{Chiffrement} &
    \tableauHeader{Ce que le fournisseur peut fournir}
}
Gmail (standard) & Serveur (Google déchiffre) &
    Contenu des emails (avec warrant), logs de connexion \\
Proton Mail & Bout en bout (E2E) &
    Métadonnées uniquement~; contenu inaccessible même avec warrant \\
iCloud (sans E2E) & Serveur (Apple déchiffre) &
    Contenu~: iCloud Backup, Photos, Notes (avec warrant) \\
iCloud (avec E2E actif) & Bout en bout (depuis iOS~16.3) &
    Métadonnées uniquement~; Apple ne peut plus déchiffrer \\
Signal & Bout en bout &
    Numéro de téléphone + date d'inscription + date dernière connexion.
    Rien d'autre~: Signal ne stocke rien. \\
\end{tableaucours}
\vspace{-0.8em}
\begin{center}{\small\itshape Tableau~17.4 --- Chiffrement E2E et disponibilité des données}\end{center}

\subsection{Résidualité des données cloud~: la mémoire longue}

À l'inverse de l'éphéméralité, certaines données cloud persistent
très longtemps dans des emplacements non évidents~:

\begin{itemize}
    \item \textbf{Sauvegardes automatiques}~: Google conserve les données
          supprimées dans ses systèmes de backup pendant 60~jours minimum
          (selon la politique de 2024)~--- non accessible à l'utilisateur,
          mais accessible via warrant.
    \item \textbf{Journaux de livraison email}~: même si un email est
          supprimé, les journaux de routage SMTP restent chez les
          intermédiaires (Google, Microsoft) pendant 30--90~jours.
    \item \textbf{Snapshots EBS abandonnés}~: des snapshots AWS créés
          par des processus de backup et non supprimés peuvent révéler
          l'état d'une instance à une date précise, souvent des mois
          après un incident.
    \item \textbf{CDN edge logs}~: CloudFront (AWS), Cloudflare, Akamai
          conservent leurs logs d'accès 30--90~jours, révélant
          des connexions à des sites même si le site lui-même a été
          supprimé.
\end{itemize}

\separateur

\begin{boxresume}
\textbf{Ce qu'il faut retenir de ce chapitre~:}
\begin{itemize}
    \item \textbf{Modèle~$\to$ preuve disponible~:} SaaS (Gmail,
          Office~365)~= logs + contenu via MLAT. IaaS (AWS EC2)~=
          logs CloudTrail uniquement~; contenu OS~= accès au compte
          client.

    \item \textbf{CloudTrail / Azure Activity Log / GCP Audit}~:
          les trois grandes plateformes IaaS logguent chaque appel API.
          Ces logs sont la source primaire de toute investigation IaaS.
          Chercher~: \texttt{StopLogging} (anti-forensique),
          \texttt{CreateUser}, \texttt{RunInstances}, \texttt{PutObject}.

    \item \textbf{Demande LEA~: 7~champs obligatoires}~: type, identifiants
          croisés, base légale (art.~52/57 Loi~2010 + Budapest/CLOUD Act),
          qualification avec double incrimination, période précise,
          distinction métadonnées/contenu, coordonnées directes.

    \item \textbf{Conservation d'abord}~: les logs de connexion
          disparaissent en 30--90~jours. Demande de conservation
          \emph{immédiate} puis MLAT dans les 60~jours.

    \item \textbf{Snapshot forensique}~: équivalent cloud de l'image
          E01. Documenter~: horodatage, ID snapshot, hash SHA-256 du
          fichier exporté, compte utilisé.

    \item \textbf{Blockchain}~: immuable, publique, pseudonyme.
          L'attribution nécessite d'identifier les exchanges KYC
          aux deux extrémités. Analyse gratuite sur blockchain.com
          ou Blockchair.

    \item \textbf{Chiffrement E2E}~: Signal/Proton~= métadonnées
          uniquement. iCloud avec E2E (iOS~16.3+)~= même Apple ne
          peut plus déchiffrer. Anticiper cette limite dans le plan
          d'investigation.

    \item \textbf{Anti-forensique cloud}~: \texttt{StopLogging} dans
          CloudTrail est lui-même loggué (s'il survient après la
          création du log)~; les logs S3 protégés par \textit{Object Lock}
          sont inviolables même par le propriétaire du compte.
\end{itemize}
\end{boxresume}

\bigskip

\begin{boxexo}[title={\ding{45}\quad Exercice~17.1 --- Modèle et preuves \niveaubase}]
Pour chacune des situations suivantes, identifiez le modèle de service
cloud impliqué, ce que le fournisseur peut fournir, et quel instrument
juridique est nécessaire~:
\begin{enumerate}[(a)]
    \item Un suspect utilise une adresse Gmail pour envoyer des emails
          de phishing. Vous voulez le contenu des emails et les IPs
          de connexion.
    \item Un ransomware a été déployé depuis une instance AWS EC2.
          Vous voulez savoir qui a lancé cette instance et depuis quelle IP.
    \item Des données de paiement volées sont stockées sur un bucket
          Dropbox. Vous voulez les métadonnées du compte et les fichiers.
    \item Un groupe WhatsApp a coordonné une fraude. WhatsApp utilise
          le chiffrement E2E. Que peut fournir Meta~?
\end{enumerate}
\end{boxexo}

\begin{boxexo}[title={\ding{45}\quad Exercice~17.2 --- Analyse CloudTrail \niveaumoyen}]
Les logs CloudTrail du compte AWS de TRANSBOIS~SA (extrait)~:
\begin{lstlisting}[language=bash_forensique]
02:10:45  GetPasswordData      IP:10.0.0.15      User:admin
02:13:45  AssumeRole           IP:185.220.XX.XX  User:admin
02:14:03  DescribeInstances    IP:185.220.XX.XX  User:admin
02:16:02  RunInstances         IP:185.220.XX.XX  User:admin
02:17:55  PutObject            IP:185.220.XX.XX  Bucket:transbois-backup
02:19:48  StopLogging          IP:185.220.XX.XX  User:admin
02:23:11  TerminateInstances   IP:185.220.XX.XX  User:admin
\end{lstlisting}
\begin{enumerate}[(a)]
    \item Identifiez les indicateurs de compromission dans cette séquence.
    \item L'action à 02h10:45 (\texttt{GetPasswordData}) précède
          la connexion suspecte à 02h13:45. Que pourrait-elle indiquer~?
    \item L'attaquant a tenté de stopper CloudTrail à 02h19:48.
          Pourquoi a-t-il échoué (hypothèse la plus probable)~?
    \item Rédigez les constatations forensiques CloudTrail et
          l'analyse correspondante avec niveau de confiance.
\end{enumerate}
\end{boxexo}

\begin{boxexo}[title={\ding{45}\quad Exercice~17.3 --- Plan d'investigation multi-cloud \niveauexpert}]
L'affaire BAOBAB implique~: fraudes MoMo vers un compte Binance~;
un compte Yahoo utilisé pour Binance~; des logs AWS suggèrent un VPN
EC2~; 3~adresses Bitcoin identifiées par analyse blockchain.
L'IP sénégalaise \texttt{41.82.XXX.XXX} apparaît dans les logs.
\begin{enumerate}[(a)]
    \item Construisez le plan d'investigation multi-cloud~: dans quel
          ordre activer quels canaux (conservation, portails LEA, MLAT,
          INTERPOL, AFRIPOL)~? Justifiez l'ordre par les durées de
          rétention.
    \item Pour chaque fournisseur (Binance, Yahoo, AWS), identifiez~:
          (i)~la base légale applicable, (ii)~le type de données
          accessible sans MLAT, (iii)~ce qui nécessite un MLAT.
    \item L'analyse blockchain révèle que les 3~adresses Bitcoin
          envoient vers un exchange turc (Paribu). La Turquie est
          partie à Budapest depuis 2014. Rédigez la demande de
          conservation urgente à envoyer via le BCN INTERPOL, en
          citant les articles applicables.
    \item Appliquez la grille CRO à la décision d'analyser publiquement
          les adresses Bitcoin (données publiques mais pseudonymes)~:
          quelles implications pour $\mathcal{C}$ si vous publiez
          l'analyse dans votre rapport judiciaire~?
\end{enumerate}
\end{boxexo}

\bigskip

\begin{boxinfo}
\textbf{Transition.} Le Chapitre~\ref{chap:anti-for} (Anti-Forensique
et Contre-Mesures) complète la Partie~IV. Il examine les techniques
utilisées pour effacer les traces dans les systèmes, les réseaux
et le cloud~--- et comment les détecter et les documenter malgré tout.
La connaissance des contre-mesures forensiques est indissociable
de la maîtrise des techniques d'investigation.
\end{boxinfo}
